% Generated draft: requires mathematical and source review.
\documentclass[11pt]{article}
\usepackage[margin=25mm]{geometry}
\usepackage{fontspec}
\setmainfont{Latin Modern Roman}
\usepackage{amsmath,amssymb,mathtools,unicode-math}
\setmathfont{Latin Modern Math}
\usepackage{enumitem,needspace,xurl,hyperref}
\setlength{\parindent}{0pt}
\setlength{\parskip}{5pt}
\setlength{\emergencystretch}{3em}
\newcommand{\N}{\mathbb{N}}
\newcommand{\Z}{\mathbb{Z}}
\newcommand{\Q}{\mathbb{Q}}
\newcommand{\R}{\mathbb{R}}
\newcommand{\C}{\mathbb{C}}
\newcommand{\F}{\mathbb{F}}
\newcommand{\A}{\mathbb{A}}
\newcommand{\PP}{\mathbb{P}}
\newcommand{\E}{\mathbb{E}}
\newcommand{\eps}{\varepsilon}
\newcommand{\dd}{\,\mathrm{d}}
\begin{document}
\section*{UnsolvedMath ID: 800809}
\textbf{Source problem code:} catalogue entry
\subsection{Definitions and mathematical statement}
Let \(M\) satisfy either of the following source hypotheses:
\begin{itemize}
\item \(M\) is a compact surface, possibly with boundary, of genus at least \(2\);
\item \(M\) is a multiply connected bounded domain in \(\R^2\) with at least two holes.
\end{itemize}
The source asks whether there exists a closed geodesic in \(\mathscr{D}^s_\mu(M)\). This notation conventionally refers to a group of volume-preserving diffeomorphisms of Sobolev regularity \(H^s\), with \(\mu\) the preserved volume measure. The supplied record does not define this group, specify \(s\) or \(\mu\), or identify the Riemannian metric determining its geodesics.

Once a Riemannian metric \(G\) and the requisite differentiable structure are specified, a nonconstant closed geodesic means a nonconstant geodesic \(\gamma\colon\R\to\mathscr{D}^s_\mu(M)\) for which there exists \(T>0\) such that
\[
\gamma(t+T)=\gamma(t)\qquad\text{for every }t\in\R.
\]
Thus the intended existence question can be recorded conditionally as follows: for each \(M\) satisfying either source hypothesis, does the intended geometric structure on \(\mathscr{D}^s_\mu(M)\) admit such a nonconstant periodic geodesic?

\emph{Unresolved qualifications.} The nonconstant interpretation is needed for a substantive existence question, but is not explicit in the source excerpt. The metric, admissible regularity, boundary conventions, and quantification over the omitted geometric data must be recovered from the source context before this becomes a fully specified mathematical statement.
\subsection{Short English statement}
For a compact surface of genus at least \(2\), possibly with boundary, or a bounded planar domain with at least two holes, the source asks whether its volume-preserving diffeomorphism group has a closed geodesic. A substantive interpretation requires a nonconstant geodesic. The supplied excerpt leaves the metric, Sobolev regularity, volume measure, and boundary conventions unspecified.
\subsection{Sources}
\begin{itemize}
\item Local UnsolvedMath registry, permanent ID \(800809\), source problem code \texttt{catalogue entry}, title: \emph{Geometric Hydrodynamics in Open Problems --- Problem 8}. The registry attributes the problem to Khesin, Misiolek, and Modin, with proposed year \(2022\), and labels it open; this label is supplied metadata, not an independently verified status assessment.
\item Khesin, Misiolek, and Modin, \emph{Geometric Hydrodynamics in Open Problems} (2022), Problem 8. Supplied source URL: \url{https://arxiv.org/abs/2205.01143}. The record reports access on 2026-10-16 and extraction by \texttt{arxiv\_latex}; these are supplied provenance details, not a claim of independent source access or verification.
\item Supplied upstream revision: \texttt{d371ee917ac9647fbb06cd6fc529deec6f3dda82}.
\end{itemize}
\end{document}
